% Compilar desde la raíz: pdflatex -output-directory=output/pdf docs/Referencias_Investigacion.tex
% Bibliografía completa de Referencias_Investigacion.bib; no requiere biber.
\documentclass[11pt,a4paper]{article}
\usepackage[utf8]{inputenc}
\usepackage[T1]{fontenc}
\usepackage{lmodern}
\usepackage[provide=*,spanish]{babel}
\usepackage[margin=2.4cm]{geometry}
\usepackage{xcolor}
\usepackage{microtype}
\usepackage{url}
\usepackage{hyperref}
\hypersetup{colorlinks=true,urlcolor=blue!55!black,linkcolor=blue!55!black,pdftitle={Referencias de la investigación: Metaheurísticas FHCP}}
\urlstyle{same}
\makeatletter
\g@addto@macro\UrlBreaks{\do\a\do\b\do\c\do\d\do\e\do\f\do\g\do\h\do\i\do\j\do\k\do\l\do\m\do\n\do\o\do\p\do\q\do\r\do\s\do\t\do\u\do\v\do\w\do\x\do\y\do\z\do\0\do\1\do\2\do\3\do\4\do\5\do\6\do\7\do\8\do\9}
\makeatother
\setlength{\parindent}{0pt}
\setlength{\emergencystretch}{5em}
\newcommand{\referencia}[1]{\par\addvspace{0.8\baselineskip}\noindent\hangindent=1.1cm\hangafter=1 #1\par\nopagebreak}
\begin{document}
\begin{center}
{\LARGE\bfseries Referencias de la investigación}\par\medskip
{\large Algoritmos heurísticos y metaheurísticas\\para el ciclo hamiltoniano (FHCP)}
\end{center}
\medskip
Bibliografía completa de las fuentes proporcionadas, organizada por área de investigación. Cada referencia incluye un enlace directo mediante DOI o URL. Las obras marcadas como segunda mano conservan esa indicación del archivo original.
\par\medskip
\textbf{Total: 88 referencias.}
\section*{Área 1. Ciclos hamiltonianos, complejidad y conjuntos de prueba}
\begin{samepage}
\referencia{\textbf{[1]} Angluin, D.; Valiant, L. G. (1979). \textit{Fast probabilistic algorithms for {Hamilton} circuits and matchings}. Journal of Computer and System Sciences, vol. 18, núm. 2, pp. 155--193.}
 {\small\url{https://doi.org/10.1016/0022-0000(79)90045-X}}\par
\end{samepage}
\begin{samepage}
\referencia{\textbf{[2]} Appel, K.; Haken, W. (1976). \textit{Every planar map is four colorable}. Bulletin of the American Mathematical Society, vol. 82, núm. 5, pp. 711--712.}
 {\small\url{https://doi.org/10.1090/S0002-9904-1976-14122-5}}\par
\end{samepage}
\begin{samepage}
\referencia{\textbf{[3]} Baniasadi, P.; Ejov, V.; Haythorpe, M.; Rossomakhine, S. (2018). \textit{A new benchmark set for traveling salesman problem and {Hamiltonian} cycle problem}. Preprint.}
 {\small\url{https://arxiv.org/abs/1806.09285}}\par
\end{samepage}
\begin{samepage}
\referencia{\textbf{[4]} Bellman, R. (1962). \textit{Dynamic programming treatment of the travelling salesman problem}. Journal of the ACM, vol. 9, núm. 1, pp. 61--63.}
 {\small\url{https://doi.org/10.1145/321105.321111}}\par
\end{samepage}
\begin{samepage}
\referencia{\textbf{[5]} Bj{\"o}rklund, A. (2014). \textit{Determinant sums for undirected {Hamiltonicity}}. SIAM Journal on Computing, vol. 43, núm. 1, pp. 280--299.}
 {\small\url{https://doi.org/10.1137/110839229}}\par
\end{samepage}
\begin{samepage}
\referencia{\textbf{[6]} Bollob{\'a}s, B.; Fenner, T. I.; Frieze, A. M. (1987). \textit{An algorithm for finding {Hamilton} paths and cycles in random graphs}. Combinatorica, vol. 7, núm. 4, pp. 327--341.}
 {\small\url{https://doi.org/10.1007/BF02579321}}\par
\end{samepage}
\begin{samepage}
\referencia{\textbf{[7]} Bondy, J. A.; Chv{\'a}tal, V. (1976). \textit{A method in graph theory}. Discrete Mathematics, vol. 15, núm. 2, pp. 111--135.}
 {\small\url{https://doi.org/10.1016/0012-365X(76)90078-9}}\par
\end{samepage}
\begin{samepage}
\referencia{\textbf{[8]} Bondy, J. A.; Murty, U. S. R. (2008). \textit{Graph theory}. Graduate Texts in Mathematics, núm. 244, Springer.}
 {\small\url{https://doi.org/10.1007/978-1-84628-970-5}}\par
\end{samepage}
\begin{samepage}
\referencia{\textbf{[9]} Chv{\'a}tal, V. (1973). \textit{Tough graphs and {Hamiltonian} circuits}. Discrete Mathematics, vol. 5, núm. 3, pp. 215--228.}
 {\small\url{https://doi.org/10.1016/0012-365X(73)90138-6}}\par
\end{samepage}
\begin{samepage}
\referencia{\textbf{[10]} Chv{\'a}tal, V.; Erd{\H{o}}s, P. (1972). \textit{A note on {Hamiltonian} circuits}. Discrete Mathematics, vol. 2, núm. 2, pp. 111--113.}
 {\small\url{https://doi.org/10.1016/0012-365X(72)90079-9}}\par
\end{samepage}
\begin{samepage}
\referencia{\textbf{[11]} Cook, S. A. (1971). \textit{The complexity of theorem-proving procedures}. Proceedings of the Third Annual {ACM} Symposium on Theory of Computing ({STOC} '71), pp. 151--158, ACM.}
 {\small\url{https://doi.org/10.1145/800157.805047}}\par
\end{samepage}
\begin{samepage}
\referencia{\textbf{[12]} Cormen, T. H.; Leiserson, C. E.; Rivest, R. L.; Stein, C. (2009). \textit{Introduction to algorithms}. 3. ed., MIT Press.}
 {\small\url{https://mitpress.mit.edu/9780262033848/introduction-to-algorithms/}}\par
\end{samepage}
\begin{samepage}
\referencia{\textbf{[13]} Dantzig, G.; Fulkerson, R.; Johnson, S. (1954). \textit{Solution of a large-scale traveling-salesman problem}. Journal of the Operations Research Society of America, vol. 2, núm. 4, pp. 393--410.}
 {\small\url{https://doi.org/10.1287/opre.2.4.393}}\par
\end{samepage}
\begin{samepage}
\referencia{\textbf{[14]} Diestel, R. (2017). \textit{Graph theory}. 5. ed., Graduate Texts in Mathematics, núm. 173, Springer.}
 {\small\url{https://doi.org/10.1007/978-3-662-53622-3}}\par
\end{samepage}
\begin{samepage}
\referencia{\textbf{[15]} Dirac, G. A. (1952). \textit{Some theorems on abstract graphs}. Proceedings of the London Mathematical Society, vol. s3-2, núm. 1, pp. 69--81.}
 {\small\url{https://doi.org/10.1112/plms/s3-2.1.69}}\par
\end{samepage}
\begin{samepage}
\referencia{\textbf{[16]} Euler, L. (1736). \textit{Solutio problematis ad geometriam situs pertinentis}. Commentarii Academiae Scientiarum Imperialis Petropolitanae, vol. 8, pp. 128--140.}
 {\small\url{https://scholarlycommons.pacific.edu/euler-works/53/}}\par
\end{samepage}
\begin{samepage}
\referencia{\textbf{[17]} Frieze, A. (2019). \textit{Hamilton cycles in random graphs: A bibliography}. Preprint.}
 {\small\url{https://arxiv.org/abs/1901.07139}}\par
\end{samepage}
\begin{samepage}
\referencia{\textbf{[18]} Frieze, A.; Haber, S. (2015). \textit{An almost linear time algorithm for finding {Hamilton} cycles in sparse random graphs with minimum degree at least three}. Random Structures \& Algorithms, vol. 47, núm. 1, pp. 73--98.}
 {\small\url{https://doi.org/10.1002/rsa.20542}}\par
\end{samepage}
\begin{samepage}
\referencia{\textbf{[19]} Garey, M. R.; Johnson, D. S.; Tarjan, R. E. (1976). \textit{The planar {Hamiltonian} circuit problem is {NP}-complete}. SIAM Journal on Computing, vol. 5, núm. 4, pp. 704--714.}
 {\small\url{https://doi.org/10.1137/0205049}}\par
\end{samepage}
\begin{samepage}
\referencia{\textbf{[20]} Gould, R. J. (2014). \textit{Recent advances on the {Hamiltonian} problem: {Survey III}}. Graphs and Combinatorics, vol. 30, núm. 1, pp. 1--46.}
 {\small\url{https://doi.org/10.1007/s00373-013-1340-x}}\par
\end{samepage}
\begin{samepage}
\referencia{\textbf{[21]} Gouyou-Beauchamps, D. (1982). \textit{The {Hamiltonian} circuit problem is polynomial for 4-connected planar graphs}. SIAM Journal on Computing, vol. 11, núm. 3, pp. 529--539.}
 {\small\url{https://doi.org/10.1137/0211042}}\par
\end{samepage}
\begin{samepage}
\referencia{\textbf{[22]} Haythorpe, M. (2018). \textit{{FHCP} challenge set: The first set of structurally difficult instances of the {Hamiltonian} cycle problem}. Bulletin of the ICA, vol. 83, pp. 98--107.}
 {\small\url{https://luca-giuzzi.unibs.it/ICA/Volumes/83//Reprints/BICA2017-32-Main-Reprint.pdf}}\par
\end{samepage}
\begin{samepage}
\referencia{\textbf{[23]} Held, M.; Karp, R. M. (1962). \textit{A dynamic programming approach to sequencing problems}. Journal of the Society for Industrial and Applied Mathematics, vol. 10, núm. 1, pp. 196--210.}
 {\small\url{https://doi.org/10.1137/0110015}}\par
\end{samepage}
\begin{samepage}
\referencia{\textbf{[24]} Helsgaun, K. (2000). \textit{An effective implementation of the {Lin--Kernighan} traveling salesman heuristic}. European Journal of Operational Research, vol. 126, núm. 1, pp. 106--130.}
 {\small\url{https://doi.org/10.1016/S0377-2217(99)00284-2}}\par
\end{samepage}
\begin{samepage}
\referencia{\textbf{[25]} Heule, M. J. H. (2021). \textit{Chinese remainder encoding for {Hamiltonian} cycles}. Theory and applications of satisfiability testing -- {SAT} 2021, Lecture Notes in Computer Science, vol. 12831, pp. 216--224, Springer.}
 {\small\url{https://par.nsf.gov/biblio/10302525-chinese-remainder-encoding-hamiltonian-cycles}}\par
\end{samepage}
\begin{samepage}
\referencia{\textbf{[26]} Hierholzer, C. (1873). \textit{{\"U}ber die {M}{\"o}glichkeit, einen {L}inienzug ohne {W}iederholung und ohne {U}nterbrechung zu umfahren}. Mathematische Annalen, vol. 6, pp. 30--32.}
 {\small\url{https://doi.org/10.1007/BF01442866}}\par
\end{samepage}
\begin{samepage}
\referencia{\textbf{[27]} Holton, D. A.; McKay, B. D. (1988). \textit{The smallest non-{Hamiltonian} 3-connected cubic planar graphs have 38 vertices}. Journal of Combinatorial Theory, Series B, vol. 45, núm. 3, pp. 305--319.}
 {\small\url{https://doi.org/10.1016/0095-8956(88)90075-5}}\par
\end{samepage}
\begin{samepage}
\referencia{\textbf{[28]} Karp, R. M. (1972). \textit{Reducibility among combinatorial problems}. Complexity of computer computations, pp. 85--103, Plenum Press.}
 {\small\url{https://doi.org/10.1007/978-1-4684-2001-2_9}}\par
\end{samepage}
\begin{samepage}
\referencia{\textbf{[29]} Koml{\'o}s, J.; Szemer{\'e}di, E. (1983). \textit{Limit distribution for the existence of a {Hamilton} cycle in a random graph}. Discrete Mathematics, vol. 43, núm. 1, pp. 55--63.}
 {\small\url{https://doi.org/10.1016/0012-365X(83)90021-3}}\par
\end{samepage}
\begin{samepage}
\referencia{\textbf{[30]} Nederlof, J. (2020). \textit{Bipartite {TSP} in $O(1.9999^n)$ time, assuming quadratic time matrix multiplication}. Proceedings of the 52nd Annual {ACM} {SIGACT} Symposium on Theory of Computing ({STOC} '20), ACM.}
 {\small\url{https://doi.org/10.1145/3357713.3384264}}\par
\end{samepage}
\begin{samepage}
\referencia{\textbf{[31]} Ohashi, R.; Soh, T.; {Le Berre}, D.; Nabeshima, H.; Banbara, M.; Inoue, K.; Tamura, N. (2025). \textit{{SAT}-based {CEGAR} method for the {Hamiltonian} cycle problem enhanced by cut-set constraints}. 28th International Conference on Theory and Applications of Satisfiability Testing ({SAT} 2025), Leibniz International Proceedings in Informatics (LIPIcs), vol. 341, pp. 24:1--24:10, Schloss Dagstuhl -- Leibniz-Zentrum f{\"u}r Informatik.}
 {\small\url{https://doi.org/10.4230/LIPIcs.SAT.2025.24}}\par
\end{samepage}
\begin{samepage}
\referencia{\textbf{[32]} Ore, {\O}. (1960). \textit{Note on {Hamilton} circuits}. The American Mathematical Monthly, vol. 67, núm. 1, pp. 55.}
 {\small\url{https://doi.org/10.2307/2308928}}\par
\end{samepage}
\begin{samepage}
\referencia{\textbf{[33]} P{\'o}sa, L. (1976). \textit{Hamiltonian circuits in random graphs}. Discrete Mathematics, vol. 14, núm. 4, pp. 359--364.}
 {\small\url{https://doi.org/10.1016/0012-365X(76)90068-6}}\par
\end{samepage}
\begin{samepage}
\referencia{\textbf{[34]} Sahni, S.; Gonzalez, T. (1976). \textit{P-complete approximation problems}. Journal of the ACM, vol. 23, núm. 3, pp. 555--565.}
 {\small\url{https://doi.org/10.1145/321958.321975}}\par
\end{samepage}
\begin{samepage}
\referencia{\textbf{[35]} Soh, T.; {Le Berre}, D.; Roussel, S.; Banbara, M.; Tamura, N. (2014). \textit{Incremental {SAT}-based method with native {Boolean} cardinality handling for the {Hamiltonian} cycle problem}. Logics in artificial intelligence -- 14th European Conference ({JELIA} 2014), pp. 684--693, Springer.}
 {\small\url{https://doi.org/10.1007/978-3-319-11558-0_52}}\par
\end{samepage}
\begin{samepage}
\referencia{\textbf{[36]} Tutte, W. T. (1946). \textit{On {Hamiltonian} circuits}. Journal of the London Mathematical Society, vol. s1-21, núm. 2, pp. 98--101.}
 {\small\url{https://doi.org/10.1112/jlms/s1-21.2.98}}\par
\end{samepage}
\begin{samepage}
\referencia{\textbf{[37]} Tutte, W. T. (1956). \textit{A theorem on planar graphs}. Transactions of the American Mathematical Society, vol. 82, pp. 99--116.}
 {\small\url{https://doi.org/10.1090/S0002-9947-1956-0081471-8}}\par
\end{samepage}
\begin{samepage}
\referencia{\textbf{[38]} Vandegriend, B.; Culberson, J. (1998). \textit{The $G_{n,m}$ phase transition is not hard for the {Hamiltonian} cycle problem}. Journal of Artificial Intelligence Research, vol. 9, pp. 219--245.}
 {\small\url{https://doi.org/10.1613/jair.512}}\par
\end{samepage}
\begin{samepage}
\referencia{\textbf{[39]} Zhou, N.-F. (2020). \textit{In pursuit of an efficient {SAT} encoding for the {Hamiltonian} cycle problem}. Principles and practice of constraint programming -- {CP} 2020, Lecture Notes in Computer Science, vol. 12333, pp. 585--602, Springer.}
 {\small\url{https://modref.github.io/papers/ModRef2019_In%20Pursuit%20of%20an%20Efficient%20SAT%20Encoding%20for%20the%20Hamiltonian%20Cycle%20Problem.pdf}}\par
\end{samepage}
\section*{Área 2. Representación del espacio de búsqueda y metaheurísticas}
\begin{samepage}
\referencia{\textbf{[40]} Ali, S.; Moscato, P. (2024). \textit{A memetic algorithm to find a {Hamiltonian} cycle in a {Hamiltonian} graph}. arXiv, Preprint.}
 {\small\url{https://arxiv.org/abs/2403.07886}}\par
\end{samepage}
\begin{samepage}
\referencia{\textbf{[41]} Ba{\c{s}}, E.; Y{\i}ld{\i}zdan, G. (2023). \textit{A discrete particle swarm algorithm with symmetry methods for discrete optimization problems}. Konya Journal of Engineering Sciences, vol. 11, núm. 3, pp. 610--634.}
 {\small\url{https://doi.org/10.36306/konjes.1199674}}\par
\end{samepage}
\begin{samepage}
\referencia{\textbf{[42]} Bean, J. C. (1994). \textit{Genetic algorithms and random keys for sequencing and optimization}. ORSA Journal on Computing, vol. 6, núm. 2, pp. 154--160.}
 {\small\url{https://doi.org/10.1287/ijoc.6.2.154}}\par
\end{samepage}
\begin{samepage}
\referencia{\textbf{[43]} Clerc, M. (2004). \textit{Discrete particle swarm optimization, illustrated by the traveling salesman problem}. New optimization techniques in engineering, Studies in Fuzziness and Soft Computing, vol. 141, Springer.}
 {\small\url{https://doi.org/10.1007/978-3-540-39930-8_8}}\par
\end{samepage}
\begin{samepage}
\referencia{\textbf{[44]} Coello Coello, C. A. (2002). \textit{Theoretical and numerical constraint-handling techniques used with evolutionary algorithms: A survey of the state of the art}. Computer Methods in Applied Mechanics and Engineering, vol. 191, pp. 1245--1287.}
 {\small\url{https://doi.org/10.1016/S0045-7825(01)00323-1}}\par
\end{samepage}
\begin{samepage}
\referencia{\textbf{[45]} Doerr, B.; Johannsen, D. (2007). \textit{Adjacency list matchings: An ideal genotype for cycle covers}. Proceedings of the Genetic and Evolutionary Computation Conference ({GECCO} '07).}
 {\small\url{https://cs.york.ac.uk/rts/docs/GECCO_2007/docs/p1203.pdf}}\par
\end{samepage}
\begin{samepage}
\referencia{\textbf{[46]} Dorigo, M.; Gambardella, L. M. (1997). \textit{Ant colony system: A cooperative learning approach to the traveling salesman problem}. IEEE Transactions on Evolutionary Computation, vol. 1, núm. 1, pp. 53--66.}
 {\small\url{https://doi.org/10.1109/4235.585892}}\par
\end{samepage}
\begin{samepage}
\referencia{\textbf{[47]} Ejov, V.; Filar, J. A.; Haythorpe, M.; Roddick, J. F.; Rossomakhine, S. (s. f.). \textit{A note on using the resistance-distance matrix to solve {Hamiltonian} cycle problem}. arXiv.}
 {\small\url{https://arxiv.org/abs/1902.10356}}\par
\end{samepage}
\begin{samepage}
\referencia{\textbf{[48]} {Flinders University} (2020). \textit{{FHCP} challenge set}. {Flinders Hamiltonian Cycle Project}.}
 {\small\url{https://sites.flinders.edu.au/flinders-hamiltonian-cycle-project/fhcp-challenge-set/}}\par
\end{samepage}
\begin{samepage}
\referencia{\textbf{[49]} Goldbarg, E. F. G.; {de Souza}, G. R.; Goldbarg, M. C. (2006). \textit{Particle swarm for the traveling salesman problem}. {EvoCOP} 2006, Lecture Notes in Computer Science, vol. 3906, pp. 99--110, Springer.}
 {\small\url{https://doi.org/10.1007/11730095_9}}\par
\end{samepage}
\begin{samepage}
\referencia{\textbf{[50]} G{\"o}rkemli, B.; Karaboga, D. (2013). \textit{Quick combinatorial artificial bee colony ({qCABC}) optimization algorithm for {TSP}}. {ISCIM} 2013, pp. 97--101.}
 {\small\url{https://core.ac.uk/download/152487671.pdf}}\par
\end{samepage}
\begin{samepage}
\referencia{\textbf{[51]} Karaboga, D.; Basturk, B. (2007). \textit{A powerful and efficient algorithm for numerical function optimization: Artificial bee colony ({ABC}) algorithm}. Journal of Global Optimization, vol. 39, núm. 3, pp. 459--471.}
 {\small\url{https://doi.org/10.1007/s10898-007-9149-x}}\par
\end{samepage}
\begin{samepage}
\referencia{\textbf{[52]} Karaboga, D.; G{\"o}rkemli, B. (2011). \textit{A combinatorial artificial bee colony algorithm for traveling salesman problem}. {INISTA} 2011, pp. 50--53.}
 {\small\url{https://avesis.erciyes.edu.tr/yayin/d8e5ca13-5345-4ab3-ab7d-27a74b551119/a-combinatorial-artificial-bee-colony-algorithm-for-traveling-salesman-problem}}\par
\end{samepage}
\begin{samepage}
\referencia{\textbf{[53]} Kirkpatrick, S.; Gelatt, Jr., C. D.; Vecchi, M. P. (1983). \textit{Optimization by simulated annealing}. Science, vol. 220, núm. 4598, pp. 671--680.}
 {\small\url{https://doi.org/10.1126/science.220.4598.671}}\par
\end{samepage}
\begin{samepage}
\referencia{\textbf{[54]} Larra{\~n}aga, P.; Kuijpers, C. M. H.; Murga, R. H.; Inza, I.; Dizdarevic, S. (1999). \textit{Genetic algorithms for the travelling salesman problem: A review of representations and operators}. Artificial Intelligence Review, vol. 13, núm. 2, pp. 129--170.}
 {\small\url{https://doi.org/10.1023/A:1006529012972}}\par
\end{samepage}
\begin{samepage}
\referencia{\textbf{[55]} Onwubolu, G. C.; Davendra, D. (2009). \textit{Differential evolution: A handbook for global permutation-based combinatorial optimization}. Studies in Computational Intelligence, núm. 175, Springer.}
 {\small\url{https://doi.org/10.1007/978-3-540-92151-6}}\par
\end{samepage}
\begin{samepage}
\referencia{\textbf{[56]} {Ponce de Le{\'o}n}, E.; Santana, R. (1999). \textit{A hybrid genetic algorithm for a {Hamiltonian} path problem}. Revista de Investigaci{\'o}n Operacional, vol. 20, núm. 1, pp. 20--29.}
 {\small\url{https://iaorifors.com/paper/29674}}\par
\end{samepage}
\begin{samepage}
\referencia{\textbf{[57]} Tasgetiren, M. F.; Liang, Y.-C.; Sevkli, M.; Gencyilmaz, G. (2007). \textit{A particle swarm optimization algorithm for makespan and total flowtime minimization in the permutation flowshop sequencing problem}. European Journal of Operational Research, vol. 177, núm. 3, pp. 1930--1947.}
 {\small\url{https://doi.org/10.1016/j.ejor.2005.12.024}}\par
\end{samepage}
\begin{samepage}
\referencia{\textbf{[58]} Thierens, D. (s. f.). \textit{Evolutionary computation}. Universiteit Utrecht.}
 {\small\url{https://ics-websites.science.uu.nl/docs/vakken/ea/slides/GArep.pdf}}\par
\end{samepage}
\begin{samepage}
\referencia{\textbf{[59]} Wagner, I. A.; Bruckstein, A. M. (1999). \textit{Hamiltonian(t): An ant-inspired heuristic for recognizing {Hamiltonian} graphs}. Proceedings of the 1999 Congress on Evolutionary Computation ({CEC} 1999), IEEE.}
 {\small\url{https://doi.org/10.1109/CEC.1999.782656}}\par
\end{samepage}
\begin{samepage}
\referencia{\textbf{[60]} Zhou, C.; Huang, L.; Wang, K.; Pang, W.; Dong, L.; Peng, L. (2003). \textit{Particle swarm optimization for traveling salesman problems}. Journal of Jilin University (Science Edition), vol. 41, núm. 4, pp. 477--480.}
 {\small\url{https://iaorifors.com/paper/46591}}\par
\end{samepage}
\begin{samepage}
\referencia{\textbf{[61]} Rothlauf, F. (2006). \textit{Representations for genetic and evolutionary algorithms}. 2. ed., Springer. \textbf{[Referencia de segunda mano.]}}
 {\small\url{https://doi.org/10.1007/3-540-32444-5}}\par
\end{samepage}
\section*{Área 4. Métricas de desempeño y evaluación experimental}
\begin{samepage}
\referencia{\textbf{[62]} Aiex, R. M.; Resende, M. G. C.; Ribeiro, C. C. (2002). \textit{Probability distribution of solution time in {GRASP}: An experimental investigation}. Journal of Heuristics, vol. 8, pp. 343--373.}
 {\small\url{https://doi.org/10.1023/A:1015061802659}}\par
\end{samepage}
\begin{samepage}
\referencia{\textbf{[63]} Aiex, R. M.; Resende, M. G. C.; Ribeiro, C. C. (2007). \textit{{TTTPLOTS}: A perl program to create time-to-target plots}. Optimization Letters, vol. 1, pp. 355--366.}
 {\small\url{https://doi.org/10.1007/s11590-006-0031-4}}\par
\end{samepage}
\begin{samepage}
\referencia{\textbf{[64]} Barr, R. S.; Golden, B. L.; Kelly, J. P.; Resende, M. G. C.; Stewart, W. R. (1995). \textit{Designing and reporting on computational experiments with heuristic methods}. Journal of Heuristics, vol. 1, núm. 1, pp. 9--32.}
 {\small\url{https://doi.org/10.1007/BF02430363}}\par
\end{samepage}
\begin{samepage}
\referencia{\textbf{[65]} Barrero, D. F.; R-Moreno, M. D.; Camacho, D. (2017). \textit{Improving experimental methods on success rates in evolutionary computation}. Journal of Experimental \& Theoretical Artificial Intelligence, vol. 29, núm. 4, pp. 695--716.}
 {\small\url{https://doi.org/10.1080/0952813X.2016.1214186}}\par
\end{samepage}
\begin{samepage}
\referencia{\textbf{[66]} Brown, L. D.; Cai, T. T.; DasGupta, A. (2001). \textit{Interval estimation for a binomial proportion}. Statistical Science, vol. 16, núm. 2, pp. 101--133.}
 {\small\url{https://doi.org/10.1214/ss/1009213286}}\par
\end{samepage}
\begin{samepage}
\referencia{\textbf{[67]} Chicco, D.; Jurman, G. (2020). \textit{The advantages of the {Matthews} correlation coefficient ({MCC}) over {F1} score and accuracy in binary classification evaluation}. BMC Genomics, vol. 21.}
 {\small\url{https://doi.org/10.1186/s12864-019-6413-7}}\par
\end{samepage}
\begin{samepage}
\referencia{\textbf{[68]} Dem{\v{s}}ar, J. (2006). \textit{Statistical comparisons of classifiers over multiple data sets}. Journal of Machine Learning Research, vol. 7, pp. 1--30.}
 {\small\url{https://www.jmlr.org/papers/v7/demsar06a.html}}\par
\end{samepage}
\begin{samepage}
\referencia{\textbf{[69]} Derrac, J.; Garc{\'i}a, S.; Molina, D.; Herrera, F. (2011). \textit{A practical tutorial on the use of nonparametric statistical tests as a methodology for comparing evolutionary and swarm intelligence algorithms}. Swarm and Evolutionary Computation, vol. 1, núm. 1, pp. 3--18.}
 {\small\url{https://doi.org/10.1016/j.swevo.2011.02.002}}\par
\end{samepage}
\begin{samepage}
\referencia{\textbf{[70]} Eiben, A. E.; Jelasity, M. (2002). \textit{A critical note on experimental research methodology in {EC}}. Proceedings of the 2002 Congress on Evolutionary Computation, {CEC} 2002, vol. 1, pp. 582--587, IEEE.}
 {\small\url{https://doi.org/10.1109/CEC.2002.1006991}}\par
\end{samepage}
\begin{samepage}
\referencia{\textbf{[71]} Hansen, N.; Auger, A.; Brockhoff, D.; Tu{\v{s}}ar, D.; Tu{\v{s}}ar, T. (2016). \textit{{COCO}: Performance assessment}.}
 {\small\url{https://arxiv.org/pdf/1605.03560}}\par
\end{samepage}
\begin{samepage}
\referencia{\textbf{[72]} Hansen, N.; Auger, A.; Ros, R.; Mersmann, O.; Tu{\v{s}}ar, T.; Brockhoff, D. (2021). \textit{{COCO}: A platform for comparing continuous optimizers in a black-box setting}. Optimization Methods and Software, vol. 36, núm. 1, pp. 114--144.}
 {\small\url{https://doi.org/10.1080/10556788.2020.1808977}}\par
\end{samepage}
\begin{samepage}
\referencia{\textbf{[73]} Hooker, J. N. (1995). \textit{Testing heuristics: We have it all wrong}. Journal of Heuristics, vol. 1, núm. 1, pp. 33--42.}
 {\small\url{https://doi.org/10.1007/BF02430364}}\par
\end{samepage}
\begin{samepage}
\referencia{\textbf{[74]} {Journal of Heuristics} (s. f.). \textit{Policies on heuristic search}. Springer.}
 {\small\url{https://www.springer.com/journal/10732/updates/17199246}}\par
\end{samepage}
\begin{samepage}
\referencia{\textbf{[75]} Rardin, R. L.; Uzsoy, R. (2001). \textit{Experimental evaluation of heuristic optimization algorithms: A tutorial}. Journal of Heuristics, vol. 7, núm. 3, pp. 261--304.}
 {\small\url{https://doi.org/10.1023/A:1011319115230}}\par
\end{samepage}
\begin{samepage}
\referencia{\textbf{[76]} Sokolova, M.; Lapalme, G. (2009). \textit{A systematic analysis of performance measures for classification tasks}. Information Processing \& Management, vol. 45, núm. 4, pp. 427--437.}
 {\small\url{https://doi.org/10.1016/j.ipm.2009.03.002}}\par
\end{samepage}
\begin{samepage}
\referencia{\textbf{[77]} Auger, A.; Hansen, N. (2005). \textit{Performance evaluation of an advanced local search evolutionary algorithm}. Proceedings of the {IEEE} Congress on Evolutionary Computation ({CEC} 2005), pp. 1777--1784, IEEE. \textbf{[Referencia de segunda mano.]}}
 {\small\url{https://doi.org/10.1109/CEC.2005.1554903}}\par
\end{samepage}
\section*{Área 5. FHCP Challenge: estándares, formatos y benchmarking}
\begin{samepage}
\referencia{\textbf{[78]} {API documentation} (s. f.). \textit{tsplib95 documentation}.}
 {\small\url{https://tsplib95.readthedocs.io/en/stable/pages/modules.html}}\par
\end{samepage}
\begin{samepage}
\referencia{\textbf{[79]} Baniasadi, P.; Ejov, V.; Filar, J. A.; Haythorpe, M.; Rossomakhine, S. (2014). \textit{Deterministic ``snakes and ladders'' heuristic for the {Hamiltonian} cycle problem}. Mathematical Programming Computation, vol. 6, núm. 1, pp. 55--75.}
 {\small\url{https://doi.org/10.1007/s12532-013-0059-2}}\par
\end{samepage}
\begin{samepage}
\referencia{\textbf{[80]} Cohen, N.; Coudert, D. (2017-12-22). \textit{Le d{\'e}fi des 1001 graphes}. Interstices.}
 {\small\url{https://interstices.info/le-defi-des-1001-graphes/}}\par
\end{samepage}
\begin{samepage}
\referencia{\textbf{[81]} {Flinders University} (2020). \textit{{FHCP} challenge set}. Conjunto de datos y soluciones. Flinders Hamiltonian Cycle Project.}
 {\small\url{http://fhcp.edu.au/fhcpcs}}\par
\end{samepage}
\begin{samepage}
\referencia{\textbf{[82]} {Flinders University} (s. f.). \textit{Flinders {Hamiltonian} Cycle Project}.}
 {\small\url{https://sites.flinders.edu.au/flinders-hamiltonian-cycle-project/}}\par
\end{samepage}
\begin{samepage}
\referencia{\textbf{[83]} Haythorpe, M.; Johnson, A. (2019). \textit{Change ringing and {Hamiltonian} cycles: The search for {Erin} and {Stedman} triples}. Electronic Journal of Graph Theory and Applications, vol. 7, núm. 1, pp. 61--75.}
 {\small\url{https://doi.org/10.5614/ejgta.2019.7.1.5}}\par
\end{samepage}
\begin{samepage}
\referencia{\textbf{[84]} Reinelt, G. (s. f.). \textit{{TSPLIB}}. Universit{\"a}t Heidelberg.}
 {\small\url{https://comopt.ifi.uni-heidelberg.de/software/TSPLIB95/}}\par
\end{samepage}
\begin{samepage}
\referencia{\textbf{[85]} Reinelt, G. (s. f.). \textit{{TSPLIB95}}. Documentaci{\'o}n en PDF; copia alojada por la Universit{\"a}t Bonn.}
 {\small\url{https://www.or.uni-bonn.de/lectures/ws17/co_exercises/programming/tsp/tsp95.pdf}}\par
\end{samepage}
\begin{samepage}
\referencia{\textbf{[86]} Applegate, D. L.; Bixby, R. E.; Chv{\'a}tal, V.; Cook, W. J. (2006). \textit{The traveling salesman problem: A computational study}. Princeton University Press. \textbf{[Referencia de segunda mano.]}}
 {\small\url{https://press.princeton.edu/books/hardcover/9780691129938/the-traveling-salesman-problem}}\par
\end{samepage}
\begin{samepage}
\referencia{\textbf{[87]} Luks, E. M. (1982). \textit{Isomorphism of graphs of bounded valence can be tested in polynomial time}. Journal of Computer and System Sciences, vol. 25, pp. 42--65. \textbf{[Referencia de segunda mano.]}}
 {\small\url{https://doi.org/10.1016/0022-0000(82)90009-5}}\par
\end{samepage}
\begin{samepage}
\referencia{\textbf{[88]} Reinelt, G. (1991). \textit{{TSPLIB}---a traveling salesman problem library}. ORSA Journal on Computing, vol. 3, núm. 4, pp. 376--384. \textbf{[Referencia de segunda mano.]}}
 {\small\url{https://doi.org/10.1287/ijoc.3.4.376}}\par
\end{samepage}
\end{document}
